% !TEX root = ms.tex
\chapter{Objetivos}
\section{Descripción del Problema}

El presente proyecto plantea la construcción teórica y práctica de una
Máquina de Cognición-Ejecución (MCE): Una arquitectura cognitiva diseñada para
abordar tareas propuestas por el usuario, planificar estrategias de resolución,
adaptarse a condiciones cambiantes del entorno y llevar a cabo acciones
computacionales que conduzcan a una solución satisfactoria.

A diferencia de los enfoques tradicionales centrados exclusivamente en el
procesamiento simbólico o en la generación de texto, la MCE se concibe como un
sistema integrado que combina razonamiento, percepción y ejecución. Su propósito
no es únicamente resolver tareas, sino hacerlo de forma autónoma, interactuando
continuamente con un entorno operativo. Esta interacción requiere no solo la
interpretación precisa de las instrucciones y la disponibilidad de acciones
programables, sino también la capacidad de adaptación reactiva frente a
imprevistos, lo cual impone exigencias particulares sobre el diseño de la
arquitectura cognitiva subyacente.

El problema central que aborda este proyecto es, por tanto, el diseño y
validación de una arquitectura cognitiva generalista que pueda ser implementada
en un sistema físico, el cual funcionará como interfaz entre el usuario, el
entorno y los módulos de decisión de la MCE. Para ello, se propone una
aproximación progresiva que comienza con el análisis teórico de los principios
fundamentales de la arquitectura, continúa con el desarrollo de un entorno
experimental para evaluar distintas configuraciones e implementaciones, y
culmina con la integración de un prototipo funcional y parcialmente completo, 
entendiendose por completo su integración desde un nivel hardware con bajos recursos
a una implementación software que satisfaga las propiedades de la MCE. 

Este prototipo, aunque limitado en términos de accesibilidad general, 
permitirá demostrar la viabilidad técnica y escalabilidad del concepto, 
con vistas a su futura aplicación en dispositivos de uso cotidiano como 
ordenadores personales o teléfonos inteligentes.

Para abordar el problema descrito, este proyecto se plantean los siguientes
objetivos tecnicos:


\begin{itemize}

	\item Definir de forma rigurosa y completa el fundamento teórico que sustenta el
    modelo de MCE propuesto.

    \item Desarrollar un entorno modular y experimental que permita probar distintas
    arquitecturas cognitivas para la implementación de la MCE.

    \item Integrar, a nivel software, las capas de Cognición y Ejecución como componentes
    diferenciados pero interoperables de la arquitectura.

    \item Implementar la MCE en un microprocesador de bajo coste, concretamente un ARM
    Cortex-A55, validando su funcionamiento en hardware embebido.

    \item Desplegar la MCE en un entorno operativo no controlado, como un sistema de
    escritorio, en el que pueda ejecutar acciones tales como abrir aplicaciones o
    navegar por internet.

    \item Habilitar la interacción entre el usuario y la MCE mediante una interfaz
    funcional, sin requerir necesariamente capacidades avanzadas de interacción
    humano-robot (HRI).

\end{itemize}


\section{Requirements}
\section{Methodology}
\section{Workplan}
